\documentclass[12pt]{article}
\usepackage[utf8]{inputenc}
\usepackage[T1]{fontenc}
\usepackage[margin=1in]{geometry}
\usepackage{setspace}
\usepackage{xcolor}
\usepackage[colorlinks=true, urlcolor=blue!60!black, linkcolor=blue!60!black]{hyperref}
\setlength{\parindent}{2em}
\setlength{\parskip}{1.25em}
\renewcommand{\baselinestretch}{1.0}
\begin{document}
\begin{center}
{\LARGE \textit{Fitzpatrick v.\ Bitzer}, 427 U.S. 445 (1976)\par}
\vspace{0.8em}
{\large Case Brief\par}
\end{center}
\vspace{1em}

\section{Memory Jogger}

Because the Fourteenth Amendment expressly limits state authority and shifted the federal--state balance, Congress may use its Section 5 enforcement power to abrogate state sovereign immunity under the Eleventh Amendment and authorize private suits for retroactive money damages and attorney's fees against unconsenting States (\href{https://www.courtlistener.com/opinion/109520/fitzpatrick-v-bitzer/}{CourtListener}; \href{https://supreme.justia.com/cases/federal/us/427/445/}{Justia}).

\section{Detailed Case Facts}

Garland M. Fitzpatrick and other present and retired male employees of the State of Connecticut brought a class action in federal court challenging sex-based disparities in Connecticut's statutory retirement benefit plan, the State Employees' Retirement Act (Conn.\ Gen.\ Stat.\ \S\S~5-162, 5-163, and 5-166) (\href{https://law.justia.com/cases/federal/district-courts/FSupp/390/278/1966559/}{District Court Opinion, 390 F.\ Supp.\ 278}; \href{https://supreme.justia.com/cases/federal/us/427/445/}{Justia}). Since a 1939 statutory amendment, Connecticut's retirement plan allowed female state employees with 25 or more years of service to retire with full pension benefits at age 50, whereas male employees with identical service had to wait five years longer, until age 55 (\href{https://law.justia.com/cases/federal/district-courts/FSupp/390/278/1966559/}{390 F.\ Supp.\ at 279--80, 283}). The plan also applied sex-based rate tables and early-retirement differentials favoring women with fewer than 25 years of service over comparably situated men (\href{https://law.justia.com/cases/federal/district-courts/FSupp/390/278/1966559/}{390 F.\ Supp.\ at 280--84}).

Originally, Title VII of the Civil Rights Act of 1964, 42 U.S.C.\ \S~2000e \textit{et seq.}, excluded state and local governments from its definition of covered ``employers'' (\href{https://supreme.justia.com/cases/federal/us/427/445/}{Justia}). On March 24, 1972, however, Congress enacted the Equal Employment Opportunity Act of 1972 (Pub.\ L.\ No.\ 92-261, 86 Stat.\ 103), explicitly relying on Section 5 of the Fourteenth Amendment to amend Title VII (\S\S~701(a), (b), and (f)) so that ``governments, governmental agencies, [and] political subdivisions'' were covered employers and civil-service workers were covered employees (\href{https://supreme.justia.com/cases/federal/us/427/445/}{Justia}). The amended statute authorized aggrieved state employees to bring private federal lawsuits against their state employers for injunctive relief, retroactive backpay (\S~706(g), 42 U.S.C.\ \S~2000e-5(g)), and ``a reasonable attorney's fee as part of the costs'' (\S~706(k), 42 U.S.C.\ \S~2000e-5(k)) (\href{https://supreme.justia.com/cases/federal/us/427/445/}{Justia}).

Sponsored by the Connecticut State Employees' Association with cooperating counsel from the Connecticut Civil Liberties Union, the male employees sued J. Frederick Bitzer (Chairman of the State Employees' Retirement Commission), the State Treasurer, and the State Comptroller (\href{https://law.justia.com/cases/federal/district-courts/FSupp/390/278/1966559/}{390 F.\ Supp.\ at 279--80}; \href{https://supreme.justia.com/cases/federal/us/427/445/}{Justia}). They sought prospective injunctive relief ending the sex-based retirement disparities, retroactive retirement benefits compensating them for past discrimination, and reasonable attorney's fees. Because the State Retirement Fund was financed predominantly by state legislative appropriations and any deficiency had to be made up from general tax revenues, all parties ultimately treated a retroactive benefit award as a money judgment payable from the Connecticut state treasury (\href{https://law.justia.com/cases/federal/district-courts/FSupp/390/278/1966559/}{390 F.\ Supp.\ at 288--89}; \href{https://supreme.justia.com/cases/federal/us/427/445/}{Justia}).

\section{Procedural History}

Plaintiffs filed suit on December 11, 1972, in the U.S. District Court for the District of Connecticut, initially raising Equal Protection and 42 U.S.C.\ \S~1983 claims before a three-judge district court; after obtaining a right-to-sue letter from the Equal Employment Opportunity Commission, they amended their complaint to add Title VII claims, and the three-judge court remanded the case to Chief Judge T. Emmet Clarie sitting alone to adjudicate the federal statutory claim first (\href{https://law.justia.com/cases/federal/district-courts/FSupp/390/278/1966559/}{390 F.\ Supp.\ at 279--81}).

On September 16, 1974, the District Court held that the Connecticut State Employees' Retirement Act violated Title VII's prohibition against sex-based employment discrimination and entered prospective injunctive relief against the respondent state officials, declining to reach the constitutional Equal Protection claim (\href{https://law.justia.com/cases/federal/district-courts/FSupp/390/278/1966559/}{390 F.\ Supp.\ at 285--90}; \href{https://supreme.justia.com/cases/federal/us/427/445/}{Justia}). However, relying on the Eleventh Amendment and \href{https://supreme.justia.com/cases/federal/us/415/651/}{\textit{Edelman v.\ Jordan}}, 415 U.S. 651 (1974), the District Court denied both retroactive retirement benefits and attorney's fees as barred monetary recoveries from the state treasury (\href{https://supreme.justia.com/cases/federal/us/427/445/}{Justia}).

The state officials did not appeal the Title VII liability holding or the prospective injunction. On the employees' appeal, the U.S. Court of Appeals for the Second Circuit (\textit{Fitzpatrick v.\ Bitzer}, 519 F.2d 559 (2d Cir.\ 1975)) agreed that Congress intended in the 1972 Amendments to authorize private backpay suits against States, but affirmed the denial of retroactive benefits, holding that under \textit{Edelman} a private federal action for retroactive damages was not a constitutionally permissible means of enforcing Fourteenth Amendment rights (\href{https://supreme.justia.com/cases/federal/us/427/445/}{Justia}; \href{https://caselaw.findlaw.com/court/us-supreme-court/427/445.html}{FindLaw}). The Second Circuit reversed and remanded as to attorney's fees, however, concluding that fee awards had only an ``ancillary effect'' on the state treasury permissible under \textit{Edelman} (\href{https://supreme.justia.com/cases/federal/us/427/445/}{Justia}).

The Supreme Court granted certiorari (423 U.S. 1031 (1975)) on both the state employees' petition (No.\ 75-251, challenging the Eleventh Amendment bar to backpay) and the state officials' cross-petition (No.\ 75-283, challenging the award of attorney's fees) (\href{https://supreme.justia.com/cases/federal/us/427/445/}{Justia}). Oral argument took place on April 20--21, 1976, and the Court issued its decision on June 28, 1976 (\href{https://supreme.justia.com/cases/federal/us/427/445/}{Justia}).

\section{Judicial Votes}

The Supreme Court ruled \textbf{9--0} in the judgment (with a 7-Justice majority opinion and 2 separate concurrences in the judgment) (\href{https://supreme.justia.com/cases/federal/us/427/445/}{Justia}):

\begin{itemize}
\item \textbf{Majority (7):} Justice William H. Rehnquist (author), joined by Chief Justice Warren E. Burger and Justices Potter Stewart, Byron R. White, Thurgood Marshall, Harry A. Blackmun, and Lewis F. Powell, Jr.
\item \textbf{Concurring in the Judgment (2):} Justice William J. Brennan, Jr.\ (separate opinion); Justice John Paul Stevens (separate opinion).
\item \textbf{Dissents (0):} None.
\end{itemize}

\section{Holding}

\textbf{Judgment:} Reversed in No.\ 75-251 (permitting retroactive backpay awards against Connecticut under Title VII) and affirmed in No.\ 75-283 (permitting attorney's fee awards against Connecticut under Title VII) (\href{https://supreme.justia.com/cases/federal/us/427/445/}{Justia}).

\textbf{Rule Announced:} The Eleventh Amendment and the principle of state sovereign immunity it embodies are necessarily limited by the enforcement provisions of Section 5 of the Fourteenth Amendment (\href{https://supreme.justia.com/cases/federal/us/427/445/}{Justia}). When Congress acts pursuant to Section 5 to enforce the substantive guarantees of the Fourteenth Amendment by ``appropriate legislation'' and clearly authorizes private suits against States, Congress may abrogate state sovereign immunity and empower federal courts to award both retroactive money damages (such as Title VII backpay) and reasonable attorney's fees payable from the state treasury (\href{https://supreme.justia.com/cases/federal/us/427/445/}{Justia}; \href{https://constitution.congress.gov/browse/essay/amdt11-6-2/ALDE_00013683/}{Constitution Annotated}).

\textbf{Holding vs.\ Dicta:} The Court's binding holdings are that: (1) Congress possesses constitutional authority under Section 5 of the Fourteenth Amendment to abrogate States' Eleventh Amendment immunity and authorize private federal suits for retroactive backpay against unconsenting States; and (2) Congress's express authorization of reasonable attorney's fees in Title VII (\S~706(k)) is likewise valid under Section 5 without needing to decide whether attorney's fees qualify as an ``ancillary effect'' under \textit{Edelman} (\href{https://supreme.justia.com/cases/federal/us/427/445/}{Justia}). Two notable points sit outside the narrow holding: first, because Connecticut did not contest that Title VII's substantive anti-discrimination rules were valid Section 5 legislation (footnote 11), the Court assumed without deciding the outer substantive boundaries of Section 5---a question later tightly policed in \href{https://supreme.justia.com/cases/federal/us/521/507/}{\textit{City of Boerne v.\ Flores}}, 521 U.S. 507 (1997); second, Rehnquist's contrast between Section 5 and Article I Commerce Clause legislation foreshadowed---without yet formally holding---the rule later adopted in \href{https://supreme.justia.com/cases/federal/us/517/44/}{\textit{Seminole Tribe of Florida v.\ Florida}}, 517 U.S. 44 (1996), that Congress cannot abrogate state sovereign immunity under Article I.

\section{Analysis of Opinions}

\subsection{Majority Opinion (Justice Rehnquist)}

Justice Rehnquist began by locating the case precisely where \href{https://supreme.justia.com/cases/federal/us/415/651/}{\textit{Edelman v.\ Jordan}}, 415 U.S. 651 (1974), had left off. In \textit{Edelman}, the Court held that the Eleventh Amendment barred retroactive welfare payments out of a state treasury because neither 42 U.S.C.\ \S~1983 nor the Social Security Act contained the ``threshold fact of congressional authorization'' to sue a State as defendant (\href{https://supreme.justia.com/cases/federal/us/427/445/}{Justia}). In \textit{Fitzpatrick}, by contrast, all parties agreed both that the plaintiffs sought retroactive monetary relief from the state treasury and that the 1972 Title VII Amendments clearly supplied congressional authorization to sue state employers (\href{https://supreme.justia.com/cases/federal/us/427/445/}{Justia}).

To resolve whether Congress could constitutionally override the Eleventh Amendment, Rehnquist focused on the structural transformation wrought by the Civil War Amendments. Ratified long after the Eleventh Amendment, the Fourteenth Amendment contains substantive prohibitions in Section 1 expressly directed at the States, paired with an affirmative enforcement grant to Congress in Section 5 (\href{https://supreme.justia.com/cases/federal/us/427/445/}{Justia}). Relying heavily on \textit{Ex parte Virginia}, 100 U.S. 339 (1880)---which described the Thirteenth and Fourteenth Amendments as ``limitations of the power of the States and enlargements of the power of Congress'' that ``carved out'' authority from state sovereignty---alongside \textit{South Carolina v.\ Katzenbach}, 383 U.S. 301 (1966), and \textit{Mitchum v.\ Foster}, 407 U.S. 225 (1972), Rehnquist concluded that the Fourteenth Amendment intentionally authorized congressional intrusions into state executive, legislative, and judicial autonomy (\href{https://supreme.justia.com/cases/federal/us/427/445/}{Justia}).

Consequently, the Eleventh Amendment and the state-sovereignty principle of \textit{Hans v.\ Louisiana} are ``necessarily limited by the enforcement provisions of \S~5 of the Fourteenth Amendment'' (\href{https://supreme.justia.com/cases/federal/us/427/445/}{Justia}). When Congress legislates under Section 5, it exercises plenary authority under an amendment specifically designed to curb state sovereignty and may therefore authorize private damages suits against States ``which are constitutionally impermissible in other contexts'' (\href{https://supreme.justia.com/cases/federal/us/427/445/}{Justia}). Finally, as to the cross-petition (No.\ 75-283), Rehnquist declined to decide whether attorney's fees were merely ``ancillary'' under \textit{Edelman}, holding instead that Congress's express authorization of fees in Title VII was directly valid under Section 5 for the exact same reason (\href{https://supreme.justia.com/cases/federal/us/427/445/}{Justia}).

\subsection{Concurring Opinion in the Judgment (Justice Brennan)}

Justice Brennan concurred only in the judgment because he rejected the majority's premise that the Eleventh Amendment even applied to suits brought by a State's own citizens (\href{https://supreme.justia.com/cases/federal/us/427/445/}{Justia}). Adhering to his dissents in \textit{Employees v.\ Missouri Public Health Department}, 411 U.S. 279 (1973), and \textit{Edelman}, Brennan maintained that the Eleventh Amendment textually bars only federal suits against States by citizens of \emph{other} States or foreign subjects. As for common-law sovereign immunity, Brennan argued that the States had surrendered that immunity ``in the plan of the Convention'' whenever they delegated enumerated powers to Congress---including both the Article I Commerce Clause and Section 5 of the Fourteenth Amendment (\href{https://supreme.justia.com/cases/federal/us/427/445/}{Justia}). Because no sovereign immunity existed to be abrogated or waived, Title VII's damages remedy was valid under both constitutional clauses.

\subsection{Concurring Opinion in the Judgment (Justice Stevens)}

Justice Stevens also concurred only in the judgment, voicing skepticism about relying on Section 5 of the Fourteenth Amendment (\href{https://supreme.justia.com/cases/federal/us/427/445/}{Justia}). In Stevens's view, Congress's Commerce Clause power amply supported regulating state employment, whereas the plaintiffs had not proved an actual Fourteenth Amendment violation and it was unclear that the 1972 Amendments were ``needed to secure the guarantees of the Fourteenth Amendment'' (\href{https://supreme.justia.com/cases/federal/us/427/445/}{Justia}). Instead, Stevens would have rejected Connecticut's Eleventh Amendment defense by narrowing \textit{Edelman}: because the backpay judgment would be paid immediately out of two ``separate and independent pension funds'' rather than directly out of general state revenues, any future state appropriations needed to replenish the trust funds were analogous to the prospective fiscal consequences permitted under \href{https://supreme.justia.com/cases/federal/us/209/123/}{\textit{Ex parte Young}} (\href{https://supreme.justia.com/cases/federal/us/427/445/}{Justia}). As to attorney's fees, Stevens would simply have treated them as ordinary litigation costs outside the Eleventh Amendment (\href{https://supreme.justia.com/cases/federal/us/427/445/}{Justia}).

\section{Examples: Future Applications}

\begin{enumerate}
\item \textbf{Same-Side (Valid Abrogation under \textit{Fitzpatrick}):} Following an extensive legislative record documenting pervasive gender stereotypes in state workplace leave policies, Congress enacts the family-care leave provision of the Family and Medical Leave Act (FMLA), explicitly authorizing state employees to sue state agencies in federal court for money damages when denied leave to care for an ailing spouse. \textit{Why it falls on the same side:} Because sex-based classifications trigger heightened Fourteenth Amendment scrutiny and Congress clearly expressed its intent to abrogate immunity to remedy documented unconstitutional state sex discrimination, the statute validly invokes Section 5 under \textit{Fitzpatrick} (as the Court later held in \textit{Nevada Department of Human Resources v.\ Hibbs}, 538 U.S. 721 (2003); see \href{https://www.law.cornell.edu/constitution-conan/amendment-14/section-5/what-may-congress-do-to-enforce-the-fourteenth-amendment-modern-doctrine}{Cornell LII Constitution Annotated}).

\item \textbf{Same-Side (Valid Abrogation under \textit{Fitzpatrick}):} Congress amends a civil rights fee-shifting statute (such as 42 U.S.C.\ \S~1988) pursuant to Section 5 to authorize federal courts to award reasonable attorney's fees against States when plaintiffs prevail on substantial Fourteenth Amendment claims, including through a favorable settlement or consent decree. \textit{Why it falls on the same side:} Directly extending \textit{Fitzpatrick}'s holding on Title VII attorney's fees, Congress may use Section 5 to authorize fee awards payable from state treasuries to incentivize enforcement of Fourteenth Amendment guarantees (as confirmed in \href{https://supreme.justia.com/cases/federal/us/448/122/}{\textit{Maher v.\ Gagne}}, 448 U.S. 122 (1980), and \href{https://constitution.congress.gov/browse/essay/amdt11-6-2/ALDE_00013683/}{Constitution Annotated}).

\item \textbf{Opposite-Side (Invalid Abrogation):} Acting solely under its Article I Indian Commerce Clause or Interstate Commerce Clause power, Congress enacts a statute expressly authorizing private parties or Indian tribes to sue unconsenting States in federal court for statutory violations. \textit{Why it falls on the opposite side:} Under \textit{Fitzpatrick}'s structural rationale---later solidified into a bright-line rule in \href{https://supreme.justia.com/cases/federal/us/517/44/}{\textit{Seminole Tribe of Florida v.\ Florida}}, 517 U.S. 44 (1996)---Article I predates the Eleventh Amendment and did not alter the original federal--state immunity balance; only a post-Eleventh Amendment grant like Section 5 of the Fourteenth Amendment (and certain narrow ``plan of the Convention'' exceptions) empowers Congress to abrogate state sovereign immunity.

\item \textbf{Opposite-Side (Invalid Abrogation):} Congress enacts a federal statute expressly subjecting States to private money-damages suits for employment discrimination based on age or disability without compiling a sufficient record of unconstitutional state conduct, sweeping far beyond what the Equal Protection Clause forbids under rational-basis review. \textit{Why it falls on the opposite side:} Although Congress invokes Section 5 and uses clear abrogation language, \textit{Fitzpatrick} applies only to ``appropriate legislation'' that genuinely enforces the Fourteenth Amendment; under \href{https://supreme.justia.com/cases/federal/us/521/507/}{\textit{City of Boerne v.\ Flores}}, 521 U.S. 507 (1997), statutes lacking ``congruence and proportionality'' between the constitutional injury and the statutory remedy exceed Section 5 and fail to abrogate immunity (see \href{https://www.naag.org/attorney-general-journal/state-sovereign-immunity/}{NAAG}; \href{https://constitution.congress.gov/browse/essay/amdt14-S5-4/ALDE_00000853/}{Constitution Annotated}).

\item \textbf{Fence-Sitter (Borderline / Contested Application):} Congress enacts a federal statute pursuant to Section 5 clearly abrogating state sovereign immunity for private damages suits against state universities that engage in disparate-impact sex discrimination (facially neutral employment policies with unequal statistical effects, even absent proof of intentional discrimination required under the Equal Protection Clause). \textit{Why it is a fence-sitter:} \textit{Fitzpatrick} upheld Title VII's abrogation without deciding whether prophylactic disparate-impact liability exceeds Section 5 (footnote 11), whereas post-\href{https://supreme.justia.com/cases/federal/us/521/507/}{\textit{City of Boerne}} doctrine demands strict ``congruence and proportionality'' to actual Fourteenth Amendment violations. Courts and scholars therefore debate whether disparate-impact claims against States remain valid prophylactic remedies for heightened-scrutiny sex discrimination or impermissibly redefine substantive constitutional rights.
\end{enumerate}

\section{Critique}

\textit{Fitzpatrick v.\ Bitzer} is a foundational pillar of modern federalism doctrine, notable in part because its majority opinion was authored by Justice Rehnquist---the chief architect of the Rehnquist Court's later expansion of state sovereign immunity (\href{https://texaslawreview.org/state-sovereign-immunity-after-the-revolution/}{Texas Law Review}). By grounding abrogation strictly in the post-Civil War shift of Section 5 rather than in Article I, Rehnquist simultaneously preserved robust federal enforcement of Reconstruction-era civil rights while laying the doctrinal foundation for \href{https://supreme.justia.com/cases/federal/us/517/44/}{\textit{Seminole Tribe of Florida v.\ Florida}}, 517 U.S. 44 (1996) (which overruled \href{https://supreme.justia.com/cases/federal/us/491/1/}{\textit{Pennsylvania v.\ Union Gas Co.}}, 491 U.S. 1 (1989), and barred Article I abrogation).

From an \textbf{originalist and textualist perspective}, \textit{Fitzpatrick} draws both praise and pointed critique. Many originalists applaud Rehnquist's ``last-in-time'' structural logic: because the Fourteenth Amendment was ratified in 1868---seventy years after the Eleventh Amendment---and explicitly restricted state sovereignty, reading Section 5 as qualifying earlier state immunity respects the constitutional text and Reconstruction's settlement (\href{https://harvardlawreview.org/print/vol-129/reconciling-state-sovereign-immunity-with-the-fourteenth-amendment/}{Harvard Law Review}). Yet stricter textualists and state-sovereignty proponents raise two objections: first, neither Section 1 nor Section 5 mentions the Eleventh Amendment or private damages suits against state treasuries, raising questions about whether the Fourteenth Amendment's ratifiers clearly intended to override state sovereign immunity without express text; second, Justice Stevens's concurrence highlighted that \textit{Fitzpatrick} brushed past whether Title VII's statutory rules actually tracked Fourteenth Amendment violations, an omission that later forced the Court in \href{https://supreme.justia.com/cases/federal/us/521/507/}{\textit{City of Boerne v.\ Flores}}, 521 U.S. 507 (1997), to impose the demanding ``congruence and proportionality'' test to prevent Congress from using \textit{Fitzpatrick} as a blank check to nullify sovereign immunity (\href{https://www.naag.org/attorney-general-journal/state-sovereign-immunity/}{NAAG}).

From a \textbf{progressive and doctrinal perspective}, scholars generally endorse \textit{Fitzpatrick}'s result but criticize the bifurcated constitutional architecture built atop it. Commentators sympathetic to Justice Brennan's concurrence argue that the Eleventh Amendment's text was never meant to bar federal-question suits by a State's own citizens in the first place, making Rehnquist's sharp wall between Article I and Section 5 historically unsupported and practically dysfunctional (\href{https://scholarship.law.unc.edu/cgi/viewcontent.cgi?article=1229&context=faculty_publications}{North Carolina Law Review}; \href{https://texaslawreview.org/state-sovereign-immunity-after-the-revolution/}{Texas Law Review}). Furthermore, progressive critics contend that once \textit{Seminole Tribe} closed off Article I abrogation and \textit{City of Boerne} narrowed Section 5 to judicially defined constitutional violations, \textit{Fitzpatrick} was transformed from a broad civil-rights shield into a narrow bottleneck---stripping state employees of damages remedies under federal age- and disability-discrimination laws while shielding States from accountability for statutory wrongs (\href{https://www.bu.edu/bulawreview/files/2014/08/REN.pdf}{Boston University Law Review}). Other scholars note an internal tension within \textit{Fitzpatrick} itself: if the Fourteenth Amendment fundamentally altered the federal--state balance, it is unclear why Section 5 legislation can pierce sovereign immunity while self-executing Section 1 guarantees (such as the Takings Clause) generally cannot without congressional action (\href{https://harvardlawreview.org/print/vol-129/reconciling-state-sovereign-immunity-with-the-fourteenth-amendment/}{Harvard Law Review}).

Ultimately, \textit{Fitzpatrick} mediates the deep tension between \textbf{Reconstruction's commitment to national enforcement of individual equality and full monetary redress against discriminatory States} and \textbf{constitutional federalism's protection of state treasuries and sovereign autonomy from unconsented private litigation}.

\section{Key Quotations}

\begin{quote}
``But we think that the Eleventh Amendment, and the principle of state sovereignty which it embodies, see \textit{Hans v.\ Louisiana}, 134 U.S.\ 1 (1890), are necessarily limited by the enforcement provisions of \S~5 of the Fourteenth Amendment.''
\end{quote}
\noindent --- Justice William H. Rehnquist, Majority Opinion, 427 U.S.\ at 456 (\href{https://www.courtlistener.com/opinion/109520/fitzpatrick-v-bitzer/}{CourtListener}; \href{https://supreme.justia.com/cases/federal/us/427/445/}{Justia}).

\begin{quote}
``When Congress acts pursuant to \S~5, not only is it exercising legislative authority that is plenary within the terms of the constitutional grant, it is exercising that authority under one section of a constitutional Amendment whose other sections, by their own terms, embody limitations on state authority. We think that Congress may, in determining what is `appropriate legislation' for the purpose of enforcing the provisions of the Fourteenth Amendment, provide for private suits against States or state officials which are constitutionally impermissible in other contexts.''
\end{quote}
\noindent --- Justice William H. Rehnquist, Majority Opinion, 427 U.S.\ at 456 (\href{https://www.courtlistener.com/opinion/109520/fitzpatrick-v-bitzer/}{CourtListener}; \href{https://supreme.justia.com/cases/federal/us/427/445/}{Justia}).

\begin{quote}
``The legislation considered in each case was grounded on the expansion of Congress' powers -- with the corresponding diminution of state sovereignty -- found to be intended by the Framers and made part of the Constitution upon the States' ratification of those Amendments, a phenomenon aptly described as a `carv[ing] out' in \textit{Ex parte Virginia}, supra, at 346.''
\end{quote}
\noindent --- Justice William H. Rehnquist, Majority Opinion, 427 U.S.\ at 455--56 (\href{https://www.courtlistener.com/opinion/109520/fitzpatrick-v-bitzer/}{CourtListener}; \href{https://supreme.justia.com/cases/federal/us/427/445/}{Justia}).

\begin{quote}
``In that circumstance, Connecticut may not invoke the Eleventh Amendment, since that Amendment bars only federal court suits against States by citizens of other States. Rather, the question is whether Connecticut may avail itself of the nonconstitutional but ancient doctrine of sovereign immunity as a bar to a claim for damages under Title VII.''
\end{quote}
\noindent --- Justice William J. Brennan, Jr., Concurring in the Judgment, 427 U.S.\ at 457 (\href{https://www.courtlistener.com/opinion/109520/fitzpatrick-v-bitzer/}{CourtListener}; \href{https://supreme.justia.com/cases/federal/us/427/445/}{Justia}).

\begin{quote}
``In my opinion the commerce power is broad enough to support federal legislation regulating the terms and conditions of state employment and, therefore, provides the necessary support for the 1972 Amendments to Title VII, even though Congress expressly relied on \S~5 of the Fourteenth Amendment.''
\end{quote}
\noindent --- Justice John Paul Stevens, Concurring in the Judgment, 427 U.S.\ at 458 (\href{https://www.courtlistener.com/opinion/109520/fitzpatrick-v-bitzer/}{CourtListener}; \href{https://supreme.justia.com/cases/federal/us/427/445/}{Justia}).


\end{document}
